\documentclass[tikz,border=6pt]{standalone}
% 공통 프리앰블 — PM Day1 교재 그림 (TikZ standalone)
\usepackage{fontspec}
\setmainfont{Apple SD Gothic Neo}
\setsansfont{Apple SD Gothic Neo}
\setmonofont{Menlo}
\usepackage{amssymb}
\usepackage{tikz}
\usetikzlibrary{arrows.meta,positioning,calc,shapes.geometric,fit,backgrounds,matrix,decorations.pathreplacing}
\definecolor{navy}{HTML}{1F3A5F}
\definecolor{teal}{HTML}{2A7F62}
\definecolor{rust}{HTML}{C8553D}
\definecolor{sand}{HTML}{E8E1D5}
\definecolor{ink}{HTML}{222222}
\definecolor{grey}{HTML}{7A7A7A}
\definecolor{lightgrey}{HTML}{EFEFEF}
\definecolor{navylight}{HTML}{DCE6F2}
\definecolor{teallight}{HTML}{DDF0E8}
\definecolor{rustlight}{HTML}{F6E0DA}
\tikzset{
  every node/.style={font=\small, text=ink},
  box/.style={draw=navy, thick, rounded corners=2pt, align=center, inner sep=5pt, fill=white},
  boxfill/.style={box, fill=navylight},
  boxteal/.style={draw=teal, thick, rounded corners=2pt, align=center, inner sep=5pt, fill=teallight},
  boxrust/.style={draw=rust, thick, rounded corners=2pt, align=center, inner sep=5pt, fill=rustlight},
  boxgrey/.style={draw=grey, thick, rounded corners=2pt, align=center, inner sep=5pt, fill=lightgrey},
  arr/.style={-{Stealth[length=2.5mm]}, thick, navy},
  arrteal/.style={-{Stealth[length=2.5mm]}, thick, teal},
  arrrust/.style={-{Stealth[length=2.5mm]}, thick, rust},
  arrgrey/.style={-{Stealth[length=2.5mm]}, thick, grey},
  lbl/.style={font=\footnotesize, text=grey, align=center},
  src/.style={font=\scriptsize, text=grey, align=left},
  title/.style={font=\normalsize\bfseries, text=navy, align=center},
}

\begin{document}
\begin{tikzpicture}[node distance=7mm]
\node[boxgrey, text width=24mm, minimum height=16mm] (in) {\textbf{입력}\\Input\\{\footnotesize 자원·예산·인력}};
\node[boxfill, text width=26mm, minimum height=16mm, right=of in] (tr) {\textbf{변환}\\Transform\\{\footnotesize 프로젝트 활동}};
\node[boxfill, text width=28mm, minimum height=16mm, right=of tr] (out) {\textbf{산출물}\\Output\\{\footnotesize 사양·일정·예산 안에서\\인도되는 것}};
\node[boxteal, text width=26mm, minimum height=16mm, right=of out] (ut) {\textbf{활용}\\Utilisation\\{\footnotesize 산출물이 실제로\\쓰이게 하는 것}};
\node[boxteal, text width=28mm, minimum height=16mm, right=of ut] (oc) {\textbf{성과 · 편익}\\Outcome\\{\footnotesize 측정 가능한 가치}};
\draw[arr] (in) -- (tr); \draw[arr] (tr) -- (out); \draw[arrteal] (out) -- (ut); \draw[arrteal] (ut) -- (oc);
% 책임 구분 브레이스
\draw[decorate, decoration={brace, amplitude=4pt, mirror}, navy, thick] ($(in.south west)+(0,-3mm)$) -- ($(out.south east)+(0,-3mm)$) node[midway, below=2mm, font=\small\bfseries, text=navy] {PM의 책임 — 산출물까지};
\draw[decorate, decoration={brace, amplitude=4pt, mirror}, teal, thick] ($(ut.south west)+(0,-3mm)$) -- ($(oc.south east)+(0,-3mm)$) node[midway, below=2mm, font=\small\bfseries, text=teal, align=center] {사업 오너(project owner)의 책임\\스폰서 또는 현업 책임자 = PRINCE2 Senior User};
\draw[rust, very thick, dashed] ($(out.east)!0.5!(ut.west)+(0,14mm)$) -- ($(out.east)!0.5!(ut.west)+(0,-22mm)$);
\node[font=\small\bfseries, text=rust, anchor=south] at ($(out.east)!0.5!(ut.west)+(0,14mm)$) {책임의 분리선};
% 온담 예시
\node[boxteal, text width=124mm, anchor=north west, font=\footnotesize, align=left] (ex) at ($(in.south west)+(0,-24mm)$) {\textbf{온담 P1의 편익 오너 (교재 3.1절)}\\
 · 폐기율 절감 43억 → 영업본부장 오정근 (구매·매장 운영 책임, Senior User) — P1 PM이 아니다\\
 · 물류 처리량 절감 96억 → 물류본부장 한동석\\
 · 매장 인건비율 38억 → 영업본부\\
 · 측정 시점 G5(2028-03-31)는 프로젝트 종료 후 → 측정 책임자는 반드시 \textbf{상설 조직의 실명}. 부서명을 쓰면 6개월 안에 유령이 된다};
\node[box, draw=rust, fill=rustlight, text width=124mm, anchor=north west, font=\footnotesize, align=left, below=3mm of ex.south west, anchor=north west] (wr) {\textbf{한국 기업에서 가장 자주 잘못 설계되는 지점} — benefit owner를 PM·PMO장으로 지정. PM은 통제 불가능한 것에 책임을 지게 되고, 범위 방어와 책임 회피가 합리적 행동이 된다. 편익 오너의 이름이 없는 편익은 착수 시점에 이미 실현되지 않기로 결정된 편익이다.};
\node[src, below=3mm of wr.south west, anchor=north west] {출처: Zwikael \& Smyrk (2011), Project Management for the Creation of Organisational Value, Springer의 ITO 모델을 바탕으로 재구성.};
\end{tikzpicture}
\end{document}
